\documentclass[11pt]{article}
\input{packages.tex}
\DeclareMathOperator{\sign}{sign}
\DeclareMathOperator{\Imm}{Im}
\DeclareMathOperator{\rg}{rg}

% Настройка колонтитулов
\pagestyle{fancy}
\fancyhf{} % Очистить все поля
\fancyhead[L]{Чил аут ребятки, копим свагу}
\fancyhead[R]{\href{https://t.me/ceagest}{Автор}}
\fancyfoot[C]{\thepage}               % Справа внизу

\setlength{\headheight}{13.6pt}

\hypersetup{
    colorlinks=true,
    linkcolor=blue,
    citecolor=blue,
    filecolor=blue,
    urlcolor=blue,
}

\begin{document}

\begin{titlepage}
    \centering
    \vspace*{\fill}  % заполняет пространство сверху

    {\Huge Лекция 2. Теорема о базисном миноре. Теорема Кронекера-Капелли. Ранг суммы и произведения матриц.
    Системы линейных уравнений. Теорема Фредгольма.
    \par}

    \vspace*{\fill}  % заполняет пространство снизу
\end{titlepage}

\setcounter{page}{1}

\section{Теорема о базисном миноре и теорема Кронекера-Капелли. Ранг суммы и произведения матриц}

\begin{theorem}[О базисном миноре]
    Любая строка матрицы является линейной комбинацией строк, содержащих базисный минор.
    Аналогичное утверждение справедливо и для столбцов.
\end{theorem}

\begin{proof}
    Докажем для строк, так как для столбцов аналогично. Рассмотрим базисный минор и строки,
    которые его содержат.
    \begin{figure}[h]
    \centering
    \includegraphics[width=0.6\linewidth]{pictures/Базисный минор.jpg}
    \end{figure} \leavevmode
    \\
    Так как в базисном миноре строки ЛНЗ, то и строки исходной матрицы, содержащие эти подстроки,
    также ЛНЗ. В самом деле, это является простым следствием доказательства теоремы о ранге подматрицы из предыдущей лекции (если бы 
    строки исходной матрицы были ЛЗ, то и строки базисного минора были бы ЛЗ, что невозможно). При этом 
    вспомним, что ранг матрицы совпадает с порядком базисного минора по определению, а значит количество 
    рассматриваемых строк совпадает с рангом матрицы, который, в свою очередь, равен строчному рангу матрицы. 
    Тогда данные строки являются базисными и, следовательно, любая строка матрицы по ним раскладывается, что и требовалось.
\end{proof}

\begin{theorem}[Кронекера-Капелли]
    Приписывание столбца $b$ к матрице $A$ не меняет её ранга тогда и только тогда, когда 
    столбец $b$ является линейной комбинацией столбцов матрицы $A$.
\end{theorem}

\begin{proof}
    $(\Longleftarrow):$ Пусть столбец $b$ есть линейная комбинация столбцов матрицы $A$.
    Припишем его к $A$ и заметим, что, коль скоро он является линейной комбинацией столбцов $A$, мы 
    можем элементарными операциями со столбцами расширенной матрицы его обнулить. Тогда вспомнив, что 
    при элементарных операциях ранг матрицы не меняется, получим:
    \[\rg \left(A \; \vdots \; b\right) = \rg \left(A \; \vdots \; 0 \right) = \rg A\]
    $(\Longrightarrow):$ Пусть теперь столбец $b$ таков, что его приписывание к матрице $A$ не меняет её ранга, 
    то есть:
    \[\rg \left(A \; \vdots \; b\right) = \rg A = r\]
    Тогда в матрице $A$ есть $r$ ЛНЗ столбцов. Посмотрим на них в расширенной матрице $\left(A \; \vdots \; b\right)$.
    В силу равенства выше и того, что эти $r$ столбцов ЛНЗ, получаем, что в расширенной матрице выбраны 
    такие ЛНЗ столбцы, что их количество равно её рангу, а значит рассматриваемые столбцы 
    являются базисными. Тогда столбец $b$ является их линейной комбинацией и, следовательно, является 
    линейной комбинацией столбцов матрицы $A$, что и требовалось.
\end{proof}

\begin{theorem}[Ранг произведения]
    Пусть $A$ и $B$ --- матрицы, для которых определено $AB$. Тогда:
    \begin{equation*}
        \begin{cases}
            \rg AB \leq \rg A \\
            \rg AB \leq \rg B
        \end{cases}
    \end{equation*}
\end{theorem}

\begin{proof}
    Для начала покажем первое неравенство. Образуем матрицу $\left(A \; \vdots \; AB \right)$ и заметим, что 
    по теореме о ранге подматрицы выполнено:
    \[\rg AB \leq \rg \left(A \; \vdots \; AB \right)\]
    Далее вспомним, как можно по другому воспринимать произведение матриц, и поймём, что 
    столбцы матрицы $AB$ являются линейной комбинацией столбцов матрицы $A$ с коэффициентами, которые берутся из матрицы $B$.
    Тогда каждый из столбцов матрицы $AB$ элементарными операциями со столбцами образованной матрицы можно обнулить. Из 
    этого прекрасного факта и того, что при ЭО ранг матрицы не меняется, получаем:
    \[\rg AB \leq \rg \left(A \; \vdots \; AB \right) = \rg \left(A \; \vdots \; O\right) = \rg A\]
    Второе неравенство можно доказать аналогично с тем только лишь изменением, что образуем следующую матрицу:
    \begin{eqnarray*}
        \begin{pmatrix}
            A \\
            \ldots \\
            AB
        \end{pmatrix}
    \end{eqnarray*}
    Также доказательство второго неравенства можно провести, опираясь на уже доказанное первое.
    Достаточно лишь транспонировать $A$ и достать из чертогов разума следующее замечательное утверждение:
    \[\rg A = \rg A^T\]
\end{proof}

\begin{definition}
    Будем говорить, что матрица $A$ имеет \textbf{упрощённый} вид, если некоторые её $r$ столбцов есть 
    столбцы единичной матрицы, а все строки, начиная с $(r + 1)$-ой, нулевые.
    \begin{figure}[h]
    \centering
    \includegraphics[width=0.54\linewidth]{pictures/Упрощённый вид.jpg}
    \end{figure}
\end{definition}

\begin{definition}
    Будем говорить, что матрица $A$ имеет \textbf{простейший} вид, если все $r$ столбцов 
    единичной матрицы стоят слева:
    \begin{figure}[h]
    \centering
    \includegraphics[width=0.45\linewidth]{pictures/Простейший вид.jpg}
    \end{figure}
\end{definition}

\paragraph{Метод Гаусса.} Пусть $A$ --- матрица. Покажем, как привести $A$ к упрощённому виду.

\begin{figure}[h]
\centering
\includegraphics[width=0.45\linewidth]{pictures/Метод Гаусса.jpg}
\end{figure}
\begin{enumerate}
    \item Находим первый ненулевой столбец слева
    \item Первое ненулевое число в нём ЭО со строками ставим в первую строку и обращаем в единицу
    \item ЭО со строками обнуляем все элементы под единицей
    \item Выполняем первые три пункта для подматрицы, образованной вычёркиванием первой строки и рассматриваемого столбца текущей матрицы
    \item Если на очередном шаге в текущей матрице нет ненулевых столбцов, то ЭО со строками обнуляем в исходной матрице все элементы, выше каждой из единиц в соответствующих столбцах, и завершаем процедуру
\end{enumerate} \leavevmode
\\
Когда дойдём до конца и получим $r$ столбцов единичной матрицы, то их количество, очевидным образом, и будет 
являться рангом $A$. В самом деле, все эти столбцы ЛНЗ, а значит ранг не менее $r$. При этом ранг не может быть больше $r$, так как ненулевых строк ровно $r$.

\begin{theorem}
    Базисную подматрицу ЭО со строками можно привести к единичной.
\end{theorem}

\begin{proof}
    Доказательство состоит в применении метода Гаусса к базисной подматрице.
\end{proof}

\begin{lemma}
    Пусть $A$ и $B$ --- матрицы равных размеров. Тогда:
    \[\rg \left(A + B\right) \leq \rg A + \rg B\]
\end{lemma}

\begin{proof}
    Рассмотрим матрицу $A + B$ и припишем к ней справа сначала матрицу $A$, а затем к полученной матрице 
    припишем матрицу $B$. Получим следующую матрицу:
    \[\left(A + B \; \vdots \; A \; \vdots \; B\right)\]
    Тогда по теореме о ранге подматрицы:
    \[\rg \left(A + B\right) \leq \rg \left(A + B \; \vdots \; A \; \vdots \; B\right)\]
    Запомним это и обозначим $\rg A = m$ и $\rg B = k$. По определению ранга это означает, что 
    в матрице $A$ есть $m$ базисных столбцов, а в матрице $B$ --- $k$ базисных столбцов. Тогда учитывая, что 
    каждый столбец матрицы $A + B$ есть линейная комбинация (даже совершенно очевидно, какая) соответствующих столбцов из $A$ и из $B$, а 
    каждый небазисный столбец каждой из матриц $A$ и $B$ раскладывается по соответствующим базисным, можем ЭО обнулить как $A + B$, так и 
    все небазисные столбцы в каждой из матриц $A$ и $B$. После сих действий в полученной матрице, ранг которой совпадает с рангом расширенной матрицы до ЭО,
    останется лишь $m + k$ ненулевых столбцов, а значит её ранг не может превзойти $m + k$. Отсюда сразу получаем требуемое неравенство:
    \[\rg \left(A + B\right) \leq \rg \left(A + B \; \vdots \; A \; \vdots \; B\right) \leq m + k = \rg A + \rg B\]
\end{proof}

\section{Системы линейных уравнений. Теорема Фредгольма}

Будем рассматривать СЛУ в самом общем виде ($m$ уравнений и $n$ неизвестных):

\begin{equation*}
    \begin{cases}
        a_{11} x_1 + a_{12} x_2 + \ldots + a_{1n} x_n = b_1 \\
        \vdots\\
        a_{m1} x_1 + a_{m2} x_2 + \ldots + a_{mn} x_n = b_m
    \end{cases}
\end{equation*} \leavevmode
\\
Пусть $A$ --- матрица коэффициентов СЛУ, $b$ --- столбец свободных членов, то есть:

\begin{equation*}
    A = 
    \begin{pmatrix}
        a_{11} & a_{12} & \ldots & a_{1n} \\
        \vdots & \vdots & \ddots & \vdots \\
        a_{m1} & a_{m2} & \ldots & a_{mn}
    \end{pmatrix}
    \quad \quad \quad
    b = 
    \begin{pmatrix}
        b_1 \\
        \vdots \\
        b_m
    \end{pmatrix}
\end{equation*} \leavevmode
\\
Тогда матрица 
\[\left(A \; \vdots \; b\right)\] 
называется расширенной матрицей системы. При этом саму систему можно 
компактно записать в виде матричного уравнения:
\[Ax = b\]

\begin{remark}
    Далее всюду при работе с СЛУ, а, если быть точным, с её расширенной матрицей, будем применять ЭО только со
    строками и лишь иногда при необходимости производить перестановку столбцов.
\end{remark}

\begin{remark}
    Так как ЭО обратимы, то, выполняя ЭО со строками расширенной матрицы СЛУ, получаем эквивалентную систему (системы эквивалентны в том смысле, что множества их решений совпадают).
\end{remark}

\begin{lemma} \label{lemma:2.1}
    Система совместна тогда и только тогда, когда 
    \[\rg \left(A \; \vdots \; b \right) = \rg A\]
\end{lemma}

\begin{proof}
    В самом деле, у системы есть решение тогда и только тогда, когда столбец её свободных членов 
    $b$ есть линейная комбинация столбцов матрицы $A$. Это утверждение, в свою очередь, по теореме Кронекера-Капелли эквивалентно 
    требуемому равенству.
\end{proof}

\begin{theorem} \label{th:2.1}
    Пусть расширенная матрица $\left(A \; \vdots \; b \right)$ СЛУ приведена к простейшему виду. Система 
    несовместна тогда и только тогда, когда в полученной матрице есть строка 
    $
        \begin{pmatrix}
            0 & 0 & \ldots & 0 & 1
        \end{pmatrix}
    $.
\end{theorem}

\begin{proof}
    $(\Longrightarrow):$ Пусть система несовместна. Тогда по теореме о ранге подматрицы и лемме \ref{lemma:2.1} имеем:
    \[\rg \left(A \; \vdots \; b \right) > \rg A = r\]
    \begin{figure}[h]
    \centering
    \includegraphics[width=0.6\linewidth]{pictures/Критерий несовместности.jpg}
    \end{figure} \leavevmode
    \\
    Но коль скоро ранг у расширенной матрицы больше, чем ранг $A$, ниже $(r + 1)$-ой строки 
    в столбце свободных членов обязательно должно стоять хотя бы одно ненулевое число, так как только приписанный 
    столбец мог продуцировать увеличение ранга. Ясно, что это число можно обратить в единицу ЭО со строками, а значит требуемое доказано.
    \\
    $(\Longleftarrow):$ Пусть в полученной матрице существует строка 
    $
        \begin{pmatrix}
            0 & 0 & \ldots & 0 & 1
        \end{pmatrix}
    $. Тогда в преобразованной системе, которая эквивалентна исходной, есть следующее равенство:
    \[0 \cdot x_1 + 0 \cdot x_2 + \ldots + 0 \cdot x_n = 1\]
    Отсюда тривиально следует, что система несовместна.
\end{proof}

\begin{definition}
    Пусть дана СЛУ, записанная в виде $Ax = b$. Транспонируем матрицу $A$ и рассмотрим новую систему:
    \[A^Ty = 0\]
    Такая система называется \textbf{сопряжённой однородной} системой для системы $Ax = b$. Иногда её 
    удобно рассматривать в следующем виде:
    \[y^T A = 0\]
\end{definition}

\begin{theorem}[Фредгольма]
    Система $Ax = b$ совместна тогда и только тогда, когда для любого решения $y$ системы $y^T A = 0$ выполнено равенство $y^T b = 0$.
\end{theorem}

\begin{proof}
    $(\Longrightarrow):$ Пусть система $Ax = b$ совместна. Тогда существует такой столбец $x$, что $Ax = b$. Значит, подействовав $y^T$ на это равенство, получим:
    \[y^T \left(A x\right) = y^T b\]
    Подставив в это равенство произвольное решение $y$ системы $y^T A = 0$, имеем:
    \[y^T b = y^T \left(A x\right) = y^T A x = \left(y^T A \right) x = 0 \cdot x = 0\]
    Но решение $y$ системы $y^T A = 0$ было выбрано произвольно, а значит требуемое доказано.
    \\
    $(\Longleftarrow):$ Пусть система $Ax = b$ несовместна. Тогда приведём её расширенную матрицу к простейшему виду и по теореме \ref{th:2.1} получим, что в этой
    матрице существует строка 
    $
        \begin{pmatrix}
            0 & 0 & \ldots & 0 & 1
        \end{pmatrix}
    $. Но в процессе преобразования расширенной матрицы выполнялись только ЭО со строками, а значит и эта строка с нулями и единицей на конце является линейной комбинацией строк расширенной матрицы системы.
    Из этого прекрасного факта сразу следует, что найдётся такой столбец $y$, что:
    \[y^T \left(A \; \vdots \; b \right) = 
    \begin{pmatrix}
        0 & 0 & \ldots & 0 & 1
    \end{pmatrix}
    \]
    Отсюда по определению умножения матриц немедленно получаем:
    \[y^T A = 
    \begin{pmatrix}
        0 & 0 & \ldots & 0
    \end{pmatrix}
    \quad \text{ и } \quad
    y^T b = 1
    \]
    Значит нашлось такое решение $y$ системы $y^T A = 0$, что $y^T b = 1 \neq 0$. Получили то, что и требовалось.
\end{proof}

\begin{definition}
    Пусть $Ax = b$ --- совместная система. Систему $Ax = 0$ назовём \textbf{приведённой} (соответствующей однородной).
\end{definition}

\begin{lemma}
    Пусть $h_1$ и $h_2$ --- решения системы $Ax = b$. Тогда $h_1 - h_2$ является решением приведённой системы $Ax = 0$.
\end{lemma}

\begin{proof}
    Так как $h_1$ и $h_2$ --- решения $Ax = b$, то:
    \begin{equation*}
        \begin{cases}
            A h_1 = b \\
            A h_2 = b
        \end{cases}
        \implies 
        A (h_1 - h_2) = 0
    \end{equation*}
\end{proof}

\begin{lemma}
    Пусть $h_1$ --- решение системы $Ax = b$, $h_2$ --- решение системы $Ax = 0$. Тогда 
    $h_1 + h_2$ является решением системы $Ax = b$.
\end{lemma}

\begin{proof}
    Так как $h_1$ --- решение системы $Ax = b$, а $h_2$ --- решение системы $Ax = 0$, то:
    \begin{equation*}
        \begin{cases}
            A h_1 = b \\
            A h_2 = 0
        \end{cases}
        \implies
        A (h_1 + h_2) = b
    \end{equation*}
\end{proof}

\begin{theorem}
    Общее решение неоднородной СЛУ есть сумма частного решения неоднородной СЛУ и общего решения однородной СЛУ.
\end{theorem}

\begin{proof}
    В самом деле, из двух лемм выше следует, что, какое бы решение $Ax = b$, отличное от частного, ни было выбрано, его 
    разность с частным решением есть решение однородной СЛУ $Ax = 0$, и при этом сумма частного решения $Ax = b$ и произвольного 
    решения $Ax = 0$ всегда является решением $Ax = b$. Это и означает, что требуемая сумма порождает все решения неоднородной СЛУ и только их.
\end{proof}

\begin{lemma}
    Пусть $Ax = 0$ --- однородная СЛУ. Если $h_1$ и $h_2$ --- решения этой системы, то $\forall \alpha, \beta \in \mathbb{R}$ число
    $\alpha h_1 + \beta h_2$ также является решением данной системы.
\end{lemma}

\begin{proof}
    Доказательство заключается в непосредственной подстановке и раскрытии скобочек по свойствам линейных операций.
\end{proof}

\begin{definition}
    Пусть $Ax = 0$ --- однородная СЛУ. Матрицу $F$ назовём \textbf{фундаментальной} матрицей для $Ax = 0$, если:
    \begin{enumerate}
        \item $AF = O$
        \item Столбцы матрицы $F$ линейно независимы
        \item Ранг матрицы $F$ максимален среди рангов всех матриц, которые удовлетворяют первым двум пунктам
    \end{enumerate}
\end{definition}

\begin{lemma}
    Пусть $A$ --- матрица, $S$ --- такая невырожденная матрица, что определено $AS$. Тогда:
    \[\rg AS = \rg A\]
\end{lemma}

\begin{proof}
    Так как матрица $S$ невырожденна, то по критерию обратимости существует $S^{-1}$. Тогда по теореме о ранге произведения матриц:
    \[\rg AS \leq \rg A = \rg \left(A S S^{-1}\right) \leq \rg AS\]
    На противоположных сторонах цепочки неравенств стоят одни и те же значения, а значит всюду должны достигаться равенства, то есть:
    \[\rg AS = \rg A\]
\end{proof}

\begin{theorem}
    Пусть $A$ --- матрица размеров $m \times n$. Пусть задана однородная СЛУ из $m$ уравнений и $n$ неизвестных $Ax = 0$. Тогда:
    \[\rg F \leq n - \rg A\]
\end{theorem}

\end{document}